\section{Evidence on the fiscal response}\label{app:fiscal}
\setcounter{table}{0}\setcounter{figure}{0}

The paper's results are reported as a function of $\hat\varphi$ and do not need its level. The reference values come from one study, however, so we bring our own evidence: Bohn-type regressions, data-determined break tests, UK fiscal events, and three designs that do not identify $\varphi$.

\paragraph{Mapping from Bohn's coefficient.} \citet{bohn1998} estimates $\rho$ in $s=\rho b+\text{controls}$. In our model the surplus responds to interest cost with slope $\varphi$, so at a steady state $ds/db=\varphi(r+\psi b)$ and $\varphi=\rho/(r+\psi b)$. Bohn's sustainability condition $\rho>(r+\psi b)-g$ is then exactly our condition $(1-\varphi)(r+\psi b)<g$. A Bohn coefficient of 0.02 at a marginal cost of 8\% corresponds to $\varphi=0.25$.

\subsection{Bohn-type regressions}\label{app:bohn}

We regress the federal primary surplus (NIPA, excluding Fed remittances, \% of GDP) on one of three lagged variables: debt held by the public, as in \citet{bohn1998}; interest payments; and the paper's consolidated debt. The controls are the output gap, temporary defense spending (defense/GDP minus its HP trend, following Barro and Bohn) and a 2020--21 dummy. The static version is Bohn's; the partial-adjustment version adds the lagged surplus and reports the long-run response. Debt coefficients are mapped to $\varphi$ with the window's mean $r+\psi b$. Because recessions both lower the surplus and raise debt, and an output-gap control removes this only in part, we repeat every regression on CBO's primary surplus with the automatic stabilizers removed (fiscal years 1967--2025, as a share of potential GDP; \citealp{cbo2026}), less Fed remittances, with CBO's GDP gap as the cyclical control.

\begin{table}[htbp]
\centering
\caption{Implied fiscal response $\varphi$ (standard error)}\label{tab:bohn}
\footnotesize
\begin{tabularx}{\textwidth}{>{\raggedright\arraybackslash}p{0.2\textwidth}>{\centering\arraybackslash}X>{\centering\arraybackslash}X>{\centering\arraybackslash}X>{\centering\arraybackslash}X>{\centering\arraybackslash}X}
\toprule
Regressor & Sample & Actual surplus, static & Actual, partial adj. & Cyclically adjusted, static & Cyclically adj., partial adj. \\
\midrule
Debt held by the public & 1967/71–2003 & 1.12 (0.51) & 1.22 (0.53) & 0.99 (0.40) & 1.39 (0.69) \\
 & 1984–2003 & $-$0.44 (0.44) & 0.63 (0.94) & $-$0.55 (0.39) & 2.61 (2.87) \\
 & 2004–2024/25 & $-$1.16 (0.09) & $-$1.23 (0.10) & $-$1.26 (0.08) & $-$1.26 (0.09) \\
 & 2004–2019 & $-$1.53 (0.17) & $-$1.71 (0.12) & $-$1.61 (0.13) & $-$1.63 (0.07) \\
Interest payments & 1967/71–2003 & 1.52 (0.51) & 2.33 (0.58) & 0.77 (0.26) & 1.62 (0.69) \\
 & 1984–2003 & $-$0.02 (0.45) & 1.03 (0.38) & $-$0.17 (0.38) & 1.72 (1.02) \\
 & 2004–2024/25 & $-$0.62 (1.41) & $-$1.95 (1.54) & 0.21 (1.67) & $-$1.63 (1.69) \\
 & 2004–2019 & 3.33 (1.25) & 2.94 (1.64) & 2.98 (1.12) & 2.49 (1.86) \\
Consolidated debt (this paper) & 1984–2003 & $-$0.48 (0.50) & 1.02 (1.26) & — & — \\
 & 2004–2024/25 & $-$1.05 (0.08) & $-$1.10 (0.09) & — & — \\
 & 2004–2019 & $-$1.42 (0.12) & $-$1.53 (0.11) & — & — \\
\bottomrule
\end{tabularx}
\par\smallskip{\footnotesize\raggedright Newey–West standard errors (2 lags); delta method for the long-run responses. Actual surplus: NIPA, calendar years, from 1971 (debt held by the public from 1970). Cyclically adjusted: CBO, fiscal years, from 1967.\par}
\end{table}


The estimates agree on direction. Every debt-based specification is lower after 2004 than before. A break imposed at 2004 is significant for debt held by the public in the static version ($p=0.005$; $p=0.002$ cyclically adjusted) and for interest payments with partial adjustment ($p=0.001$ in both). It is not significant for the consolidated debt measure ($p\approx0.6$), whose pre-2004 window starts only in 1981 and is itself negative in the static version. The one exception is the interest specification over 2004--19, a period when interest costs were low and falling while the 2011 spending caps cut deficits; its standard error exceeds 1.

They do not agree on size, even without the automatic stabilizers. Implied $\varphi$ ranges from about $-2$ to $+3$ across samples and specifications, far outside any plausible structural value. Removing the stabilizers changes the estimates little, so the business cycle is not the main problem: the surplus also moves with revenue booms, tax legislation and wars that are not responses to debt. That is why \citet{auerbach2024} use legislated changes, and we do not replace their estimates with ours.

The results that matter do not need the level. The inflation layer, and with it the effect of QE, does not depend on $\hat\varphi$ at all. The ranking of 2023--25 as the largest requirement since 1980 survives without assuming a collapse of the response at any date: with one common $\hat\varphi$ in every year, it holds for any $\hat\varphi$ below about 0.35 (Figure~\ref{fig:limitmap}, line). At $\hat\varphi=0.25$ throughout, the 2023--25 requirement is 0.46--0.53pp a year, against a pre-2022 maximum of 0.43 in 1981. Only if the pre-2004 legislated value of 0.39 held in every year does 1981 rank first.

\subsection{Where does the response break?}\label{app:breaks}

Imposing 2004 assumes the answer. We let the data choose the break dates \citep{bai1998,bai2003}: all coefficients of the static regression may change at each break, the minimum segment is 15\% of the sample, 2020--21 are dropped, and $p$-values come from a wild bootstrap.

\begin{table}[htbp]
\centering
\caption{Data-determined breaks in the U.S. fiscal response}\label{tab:breaks}
\small
\begin{tabularx}{\textwidth}{>{\raggedright\arraybackslash}p{0.3\textwidth}>{\centering\arraybackslash}X>{\centering\arraybackslash}X}
\toprule
Surplus, regressor & Breaks: one / two / three & Number of breaks (sequential / BIC) \\
\midrule
Actual (NIPA, 1971–2024), debt & 1992 / 1992, 2013 / 1979, 1992, 2013 & 1 / 3 \\
Actual, interest & 2009 / 1992, 2009 / 1983, 1995, 2012 & 0 / 3 \\
Cyclically adjusted (CBO, 1967–2024), debt & 1993 / 1979, 1993 / 1979, 1996, 2012 & 1 / 3 \\
Cyclically adjusted, interest & 2009 / 1996, 2009 / 1979, 1996, 2009 & 3 / 3 \\
\bottomrule
\end{tabularx}
\par\smallskip{\footnotesize\raggedright \emph{Notes.} The test of no break against the best number of breaks (UDmax) has a bootstrap p-value of 0.00 in every row.\par}
\end{table}


The response did change, but the data place the changes in the mid-1990s and around 2009--13, not at 2004. The segment slopes are not credible responses: after 1992 the debt slope is negative, because the surpluses of the 1990s reduced debt, and before 1991 the interest slope exceeds 2. Time-series regressions across regimes can date the changes in behaviour but cannot measure them.

\subsection{UK fiscal events}\label{app:obr}

A cleaner design follows \citet{auerbach2024}: at each fiscal event, compare the policy decisions with the revision to projected debt interest that the government saw before deciding. The UK's Office for Budget Responsibility publishes this decomposition for every event since 2010. Its fiscal forecast revisions database splits each revision to the borrowing forecast into policy decisions and underlying changes, and splits the underlying changes into receipts, debt interest and other spending. Debt interest (net of the Asset Purchase Facility) moves with gilt yields, Bank Rate and, through index-linked gilts, RPI inflation, none of which the Chancellor controls. For each of 32 events from November 2010 to March 2026 we average over the five years after the current one, in \% of GDP, and regress the policy change on the debt-interest revision, controlling for the receipts and non-interest revisions; $\varphi$ is minus the slope. To remove the part of the debt-interest revision that comes from revisions to borrowing, we instrument it with the revisions between consecutive forecasts to the market assumptions behind it: short rates and gilt rates, and the RPI price level.

\begin{table}[htbp]
\centering
\caption{UK: policy response to revisions in projected debt interest, 2010–2026}\label{tab:obr}
\small
\begin{tabularx}{\textwidth}{>{\raggedright\arraybackslash}p{0.3\textwidth}>{\centering\arraybackslash}X>{\centering\arraybackslash}X}
\toprule
Specification & $\varphi$ (s.e.) & First-stage F \\
\midrule
OLS, all 32 events & 0.30 (0.19) & — \\
OLS, excluding the COVID events (29) & 0.19 (0.10) & — \\
IV, rate and RPI assumption revisions, all events & 0.00 (0.31) & 20 \\
IV, excluding the COVID events & 0.15 (0.10) & 66 \\
\bottomrule
\end{tabularx}
\par\smallskip{\footnotesize\raggedright \emph{Notes.} HC1 standard errors. COVID events: March 2020, November 2020, March 2021. Dropping one event at a time moves the OLS estimate between 0.09 (without November 2022) and 0.41.\par}
\end{table}


When higher debt interest comes from the market, UK policy has offset about 0.15 of it, mostly through spending, under explicit fiscal rules. The estimate is imprecise and leans on November 2022, when gilt yields and RPI inflation rose together and the Autumn Statement consolidated. It supports a low $\hat\varphi$, within the paper's range of 0 to 0.25, but the results do not depend on it.

\subsection{\texorpdfstring{Designs that do not identify $\varphi$}{Designs that do not identify phi}}\label{app:fail}

\paragraph{The clock as an instrument, U.S. annual data (1980--2024).} The predicted interest-cost change $b_t P_t(1)\,\Delta r_{t+1}$ uses the predetermined repricing exposure as an instrument for interest-cost changes; the outcome is the subsequent change in the NIPA primary surplus, excluding Fed remittances, and the controls are the rate change, debt, growth and the lagged surplus. The first-stage $F$ is about 6.7 and falls below 1 once $\Delta r\cdot b$ is controlled, and the implied $\varphi$ is implausible (6--26). The U.S. repricing speed varies too little over time: $P(1)$ has a standard deviation of 0.056, and the instrument is 99\% correlated with $\Delta r\cdot b$.

\paragraph{A cross-country panel (21--23 advanced economies).} The same instrument is built with two measures of the share of debt that reprices within a year: the OECD Economic Outlook's refinancing share (2014--2025) and the BIS bills share (1990--2025). The outcome is the underlying or actual primary balance over local-projection horizons of 0--3 years, with country and year fixed effects. With the refinancing share, the first stage is adequate ($F=9$--15) but $\varphi$ comes out at 1.5--3.2. The usable variation is $\Delta r\cdot b$, which is confounded: in 2022--24 the countries with the most debt and the largest rate increases were also unwinding the most COVID and energy support. Once $\Delta r\cdot b$ is controlled, the instrument is weak ($F=1$--5). The BIS bills share is inconsistent across countries: some shares exceed 1, and the totals and the short-term series do not always cover the same instruments. The design is close to \citet{andreolli2024}, who interacts debt duration with monetary policy shocks, and to \citet{johns2026}, who interact high-frequency euro-area monetary shocks with the level and maturity of public debt. Neither estimates the offset $\varphi$, but \citet{johns2026} find that primary balances deteriorate after monetary tightening, which is consistent with a weak fiscal response.

\paragraph{U.S. CBO baseline updates (1992--2026).} This is the same design as for the UK, using CBO's record of changes between baselines by source (legislative, economic and technical) from its public evaluation repository. Policy is the legislative change in the primary deficit, the interest revision is the economic plus technical change in net interest, the controls are the non-legislative changes in revenue and non-interest outlays, and the instruments are changes in ten-year and three-month yields between updates. The estimates are uninformative: OLS gives 0.28 (0.51), the IV first stage is weak ($F\approx5$), and lagged responses over the next one to four updates lie between $-1.5$ and 0.5, with standard errors of 0.5--0.9. Legislative changes are dominated by a few large packages (the 2001--03 tax cuts, the 2009 Recovery Act, the 2017 tax act and the 2020--21 COVID relief), whose standard deviation, 0.55\% of GDP, is three times that of the interest revisions (0.18\%).
